\section{核磁共振}
核磁共振\textrm{(Nuclear Magnetic Resonance, NMR)}谱是研究分子、液体和固体的化学环境的高灵敏技术。通过第一原理量子力学可以模拟和解析\textrm{NMR}谱，即计算化学屏蔽、核四级耦合常数和超精细耦合常数。

\subsection{\rm{NMR}:~化学屏蔽}
\subsubsection{金刚石中的汇聚化学屏蔽}
对金刚石的双原子初基原胞计算化学屏蔽，检查其随着平面波截断能和$\vec k$点的收敛情况。

在磁场中，有非零自旋$\mathbf{I}$的原子核，当外磁场为$\mathbf{B}_{\mathrm{ext}}$，进动磁矩为$\mathbf{\mu}$(参见下图中双倒反圆锥上的红箭头)。该磁矩的进动的\textrm{Larmor}频率为$\omega_L$，其大小由外磁场$\mathbf{B}_{\mathrm{ext}}$和原子核的磁-旋比$\gamma$的\textcolor{red}{相互作用强度}。如果振动磁场，即当施加的射频脉冲\textrm{(Radio-Frequency (RF) Pulse)}的频率为$\omega_{\mathrm{rf}}$，并且垂直\textcolor{red}{(横向平面)}恒定外磁场$\mathbf{B}$(参考面)时，核磁矩发生改变。如果射频频率$\omega_{\mathrm{rf}}$与$\omega_L$接近，就会发生磁共振。换言之，核磁矩会由平行于磁场变成垂直磁场，即指向横向平面。然后，通过磁矩的弛豫，回到参考平面，产生核磁共振中测量的信号。

看起来相同的核同位素，即使具有相同的旋磁比$\gamma$，也可能会因化学环境不同而具有不同的共振频率。围绕在核外的不同壳层的电子会对核产生屏蔽。例如，在溴乙烷(\ch{Br-CH2-CH3})中，在空间上，亚甲基\ch{-CH2-}官能团的\ch{1H}比起甲基\ch{-CH3}更接近\ch{Br}。化学屏蔽改变了原子核处磁场，因此出现的核磁共振频率也发生了变化。相对标准参考发生的频率移动称为化学移动$\delta$。可以用\textrm{NMR}来作为探测分子和固体的化学结构。

核磁共振实验中测量的化学位移$\delta$，标准参考是四甲基硅烷\textrm{(Tetramethyksilane, TMS)}，参考化合物的化学屏蔽是$\sigma^{\mathrm{ref}}$，如果样品的化学屏蔽是$\sigma$，则有
\begin{displaymath}
	\delta_{ij}(\mathbf{R})=\sigma_{ij}^{\mathrm{ref}}=\sigma_{ij}^{\mathrm{ref}}-\sigma_{ij}(\nathbf{R})
\end{displaymath}
这里$ij$遍历全部\textrm{Cartesian}方向，$\mathbf{R}$是原子核位置。

根据线性响应，化学屏蔽张量可以通过下式表示
\begin{displaymath}
	\sigma_{ij}(\mathbf{R})=-\dfrac{\partial B_i^{\mathrm{ind}}(\mathbf{R})}{\partial B_j^{\mathrm{ext}}}
\end{displaymath}
这里$B^{\mathrm{inc}}$是诱导磁场。

本算例中，计算金刚石结构的\ch{C}元素的各向同性的化学屏蔽值随着平面波截断能和$\vec k$点网格的收敛情况。

计算所需要的输入文件在\textrm{shielding}目录下。

金刚石的初基原胞结构文件\textrm{POSCAR}

计算参数控制文件\textrm{INCAR.nmr}

$\vec k$空间网格剖分文件\textrm{KPOINTS.nmr}

采用$4\times4\times4$的网格剖分精度。

检查\textrm{INCAR}文件中的控制计算参数。参数\texttt{ENCUT}、\texttt{ISMEAR}、\texttt{SIGMA}、\texttt{EDIFF}、\texttt{PREC}和\texttt{LASPH}就是一般的\textrm{DFT}计算的设置。需要说明的是，对于\textrm{NMR}计算，电子态结构计算的精度要比较高(\texttt{PREC}=\textrm{A}，\texttt{EDIFF}$\glantleq$1E-8和高的\texttt{ENCUT})。

其余的计算参数用于化学屏蔽的计算。\texttt{LCHIMAG}表示通过线性响应计算化学屏蔽。\texttt{LNMR\_SYM\_RED}表示当采用线性响应方法计算化学屏蔽时，放弃与$\vec k$空间求导不一致的对称操作。\texttt{NLSPLINE}要求\textrm{PAW}方法的投影函数在倒空间是可微分的。\texttt{KPAR}用于极速计算，确定的是\textrm{VASP}并行计算中每一次同时处理的$\\vec k$点数目。

首先检查计算化学屏蔽随平面波截断能的收敛情收敛情况。为了在不同截断能下执行\textrm{VASP}计算，\textrm{INCAR}需要作相应调整。采用\textrm{Bash}脚本\textrm{encut\_shielding.sh}，分别用\textrm{500-900~eV}间，步长依次为\textrm{100~eV}的值依次替代原来\textrm{INCAR}文件中的\texttt{ENCUT}=\textrm{400~eV}。在执行目录下，执行该脚本。

计算完毕后，脚本还会通过\textrm{Bash}命令
\fbox{\textrm{grep -A 5``iso\_shield'' OUTCAR | tail -1 | awk '\{print \$6\}'}}
就可以在\textrm{OUTCAR}文件中提取\ch{13C}的化学屏蔽值。

在计算过程中，标准输出在不同阶段的输出，首先是\textrm{DFT}计算的电荷\textrm{(CHGCAR)}和波函数\textrm{(WAVECAR)}的收敛。随后沿着磁场轴向线性响应的输出，\textsc{BDIR}表示外加磁场，\textsc{IDIR}表示诱导磁场。数值表示\textrm{Cartesian}轴向(1:$x$-轴；2-$y$-轴；$3$-轴)。\textsc{IQ}是有限差分下的格点的值。格点的定义由计算控制参数\texttt{ICHIBARE}定义，取值范围是-\texttt{ICHIBARE}到+\texttt{ICHIBARE}。在本算例中，因为\texttt{ICHIBARE}取的是默认值1，因此有3个\texttt{IQ}点。

\textrm{Bash}脚本\textrm{encut\_shielding.sh}执行过程中，不同能量截断下执行\textrm{VASP}计算化学屏蔽的数据保存在文件\textrm{encut\_shielding.dat}。在\textrm{OUTCAR}文件中，关于化学屏蔽的输出为

\begin{table}[h!]
    \centering
    \begin{tabular*}{\textwidth}{@{\extracolsep{\fill}}c@{\extracolsep{\fill}}
        \toprule
	名称 &含义\\
\midrule
\texttt{EXCLUDING G=0 CONTRIBUTION} & 排除宏观长程贡献的化学屏蔽，该参数对于分子很有用\\
\texttt{INCLUDING G=0 CONTRIBUTION} & 包含宏观长程贡献的化学屏蔽，特别是对于表面电流，该参数对于晶体很有用\\
\texttt{ion} & \textrm{POSCAR}文件中排序的原子数目\\
\texttt{BDIR} & 施加磁场的方向(即\textsc{B0})，比如，这里\texttt{BDIR}=1 表示沿$x$轴方向。每一列的外加磁场可以分解为$x$，$y$，$z$方向的贡献，所以对每个原子来说，这才是整个化学屏蔽张量 \\
\texttt{iso\_shield} & 化学屏蔽张量之后，括号里给出各向同性的化学屏蔽，跨度\textrm{(Span)}和斜度\textrm{(Skew)}。各向同性的化学屏蔽，即化学屏蔽张量的对角元平均值($\sigma_{\mathrm{iso}}=\mathrm{Tr}(\sigma)/3=(\sigma_{11}+\sigma_{22}+\sigma_{33})/3$) \upcite{SSNMR2-28590010_1993}\\
\texttt{span} & 跨度\textrm{(Span)}$\Omega$定义为 $\Omega=\sigma_{33}-\sigma_{22}$ \upcite{SSNMR2-28590010_1993}\\
\texttt{skew} & 斜度\textrm{(Skew)}$\kappa$定义为 $\kappa=3(\sigma_{\mathrm{iso}}-\sigma_{22})/(\sigma_{33}-\sigma_{11})=3(\sigma_{\mathrm{iso}}-\sigma_{22})/\Omega$ \upcite{SSNMR2-28590010_1993}
        \bottomrule
    \end{tabular*}
    \caption{\textrm{OUTCAR}文件中关于\textrm{NMR}的参数}}
\label{Tab:-Tag}
\end{table}

注意这里的化学屏蔽包含了主要的贡献，这些是最后与实验对比的重要的数据项。这些数据涵盖了下一个算例的内容。

通过\textrm{Py4vasp}绘制每个能量截断下的化学屏蔽

截断能达到\textrm{500~eV}以上，总能量收敛会达到\textrm{10~meV}。如果截断能低于该值，对波函数来说，平面波数目太少，不再是好的展开基组。对于化学屏蔽来说，达到收敛的截断能会需要取得很高，截断能达到\textrm{500~eV}，化学屏蔽快速收敛到\textrm{1~ppm}，但是，如果要将化学屏蔽收敛到\textrm{0.1 ppm}，截断能至少要\textrm{800~eV}。

注意，计算所能达到的精度与元素也有关系。例如，对于\ch{1H}来说，典型的化学屏蔽范围是\textrm{0-14~ppm}，而对于\ch{13C}，化学屏蔽的范围是\textrm{0-250~ppnm}。这些值并不与化学屏蔽严格对应，但给出了数量级上的参考。化学屏蔽可以给出到\textrm{0.1~ppm}的精度。对于重元素，化学屏蔽的范围更大，比如\ch{195Pt}，范围将会达到-7000到\textrm{+7000~ppm}，这种情况下，\textrm{0.1~ppm}不是一个可行的精度要求，必须另外选取精度(如\textrm{1~ppm})。

一般地，这里并不推荐计算参数随总能的收敛。
